\section{An elementary upper bound for corners}\label{sec:elementary}
We prove that every fifteen sources have a corner assembly. The argument
uses the two kinds of repeated colour differently: opposite-repeat cubes
supply both orientations of certain corners, whereas adjacent-repeat cubes
supply one orientation and two corners involving their repeated colour.
The longer geometric reductions are proved in the separately paginated
\emph{Mathematical appendix to Sharp thresholds for five-colour cube
assemblies}. That appendix is part of the mathematical argument, intended
for refereeing with this article; the checking programs are supplementary.

\subsection{Directed corners and repeated colours}
Encode an oriented tricoloured corner $(i,j,k)$ by the complementary ordered
pair $ab$, chosen so that $(a,b,i,j,k)$ is an even permutation of the five
colours. Cyclic permutation of $(i,j,k)$ does not change $ab$; reversing it
replaces $ab$ by $ba$. Write $V=\col$.
The six tricoloured corners of an adjacent-repeat variety form a directed cut
\[
 C(S)=\{ab:a\in S,\ b\in V\setminus S\},\qquad |S|\in\{2,3\}.
\]
The smaller side $P$ of the cut is its \emph{pole pair}. For each cut there
are three varieties, indexed by their repeated colour $r\notin P$.
Their other two corners are $rrp$, for the two $p\in P$.
Thus there are twenty adjacent classes, each consisting of three varieties.
Class multiplicity always counts physical copies, allowing different
repeated colours within a class.
An opposite-repeat variety is represented by a four-cycle on four vertices
of $V$: it supplies both directed corners on each edge, and the omitted
vertex is its repeated colour. Equivalently its edges form a $K_{2,2}$.
These descriptions follow by reading the corners around the repeated faces.

For a class with at least two copies, call a six-corner matching after
reserving two of its copies an \emph{own completion}. The reserved copies
need not have the same repeat, so this alone is not an assembly.

\begin{lemma}[Repeated-colour completion]\label{lem:repeat-completion}
An adjacent class with an own completion assembles its target if two of
its copies have the same repeated colour. If two classes each have three
copies and an own completion, and their pole pairs intersect, the pool
has an assembly.
\end{lemma}
\begin{proof}
A same-repeat pair fills the two bicoloured corners, leaving precisely the
six-corner matching. Otherwise each triple must use all three repeats
allowed by its pole pair. The two repeat sets intersect; choose a common
colour $r$. Reserve one copy from each class, using the second class at
the shared pole of the first target. Start with the first class's own
completion. Restore one of its two reserved copies and remove the chosen
copy of the second class. If the latter was used, the restored copy,
which fits every tricoloured corner of this target, replaces it. The
remaining six positions are still matched to distinct sources.
\end{proof}

\begin{lemma}[A nine-source construction]\label{lem:nine-core}
For distinct colours $a,b,c,d,e$, an opposite source containing the edge
$ab$, together with three sources in each of $C(ac),C(bc)$ and two in
$C(abc)$, has an assembly. The same holds after reversing all cuts.
\end{lemma}
\begin{proof}
Put $A=C(ac)$, $B=C(bc)$ and $H=C(abc)$. After reserving two $A$ copies,
match its six tricoloured corners by
\[
 A:cb,\qquad B:cd,ce,\qquad H:ad,ae,\qquad O:ab.
\]
After reserving two $B$ copies use
\[
 B:ca,\qquad A:cd,ce,\qquad H:bd,be,\qquad O:ba.
\]
A class assigned two positions supplies two distinct copies. Both triples
have own completions, and their poles meet at $c$, so
Lemma~\ref{lem:repeat-completion} applies.
\end{proof}

\subsection{Structural inputs}
The appendix proves the following concentration and assembly statements.
Their proofs use the cut model, Hall's theorem, repeated-colour completion,
and explicit geometric cases. None assumes the corner threshold.
The prescribed-family lemma~\ref{lem:eight-common-missing}, proved earlier,
is used in two concentrated-opposite cases.

\begin{lemma}[Concentration]\label{lem:elementary-concentration}
A failing fourteen-source pool containing at least three copies of an
opposite variety has one of these two profiles: four copies of that variety
and adjacent-class sizes $5,5$, or three copies and adjacent-class sizes
$5,3,3$. In either case there are no other opposite sources.
A mixed fourteen-source pool in which every opposite variety has
multiplicity at most two and some adjacent class has multiplicity at least
four has an assembly.
\end{lemma}
The opposite profiles are proved in the appendix arguments headed
``Four identical opposite copies'' and ``Three identical opposite copies'';
the mixed concentration arguments then treat the possible opposite totals.
The conclusion records class sizes, not repeated-colour refinements.

\begin{lemma}[Assembly inputs]\label{lem:elementary-inputs}
The following pools have assemblies:
\begin{enumerate}
\item fifteen adjacent sources;
\item one opposite source and fourteen adjacent sources;
\item three opposite sources and eleven adjacent sources, with at most
three copies in each adjacent class;
\item five opposite sources, each variety occurring at most twice,
and sources $2A+B+C$ in three distinct adjacent classes;
\item nine adjacent sources in distinct classes;
\item eight opposite sources, each variety occurring at most twice;
\item seven opposite sources with that same cap and any adjacent source.
\end{enumerate}
\end{lemma}
The appendix gives complete proofs, including the geometric exhaustions
behind the finite Hall tables. Item~3 includes all three opposite
multiplicity patterns, including three identical varieties.
Items~4--7 produce opposite targets; the other items may produce either
kind of target. These distinctions will matter when a source is replaced.

For the present consolidation draft, the principal appendix locations are:
% Generated by assemble_appendix.py; do not renumber by hand.
\[\begin{array}{l|l}
\text{Input}&\text{Appendix argument}\\\hline
\text{Concentrated opposite profiles}&22,14\\
\text{Fifteen adjacent sources}&39\\
\text{One opposite and fourteen adjacent}&52\\
\text{Three opposite and eleven adjacent}&66\\
\text{Five-opposite helper construction}&30\\
\text{Nine distinct adjacent classes}&29\\
\text{Eight opposite; seven plus one adjacent}&10\\
\text{Butterfly; two stars; star and triangle}&54,63,59\\
\end{array}\]


These locations include references to the supporting geometric arguments;
they do not replace those arguments.

For orientation, the proof of item~1 is a deletion argument. A failing
pure fourteen-pool has class cap three, at most eight occupied classes,
and at most five repeated classes. If none of its triples has an own
completion, its class sizes are $3,3,3,3,1,1$. A failing pure fifteen-pool
with class cap three cannot have two own-completed triples. Deleting a
copy of any triple leaves a doubled class, hence forces an own-completed
triple elsewhere. Restore the copy and delete from this new triple:
a second own-completed triple is forced, a contradiction.
For item~2, the appendix proves that a pure fourteen-failure has at least
two triples. A triple without an own completion has a private column in
its two-by-three corner array. The added opposite source either meets
that column, creating a second own completion, or belongs to one of five
avoiding varieties, whose three geometric types have direct completions.
These explanations are reading guides; the appendix supplies the required
concentration, private-column and repeated-colour proofs.

\subsection{Compressing two adjacent sources}
\begin{lemma}[Disjoint-cut compression]\label{lem:compression}
Two disjoint adjacent cuts have an opposite variety's eight corners in
their union. Replacing two actual sources from these classes by one
virtual source of that opposite variety preserves the implication
``the smaller pool assembles $\Rightarrow$ the original pool assembles''.
\end{lemma}
\begin{proof}
Normalize the first cut to $C(01)$. The disjoint cuts are exactly
$C(234),C(23),C(24),C(34)$. For $C(ij)$ take the opposite variety
with bipartition $01\mid ij$; for $C(234)$ choose any two of $2,3,4$.
The first actual source supplies one directed half and the second the
other. An assembly assigns only one position to the virtual source.
Replace it by whichever removed source supplies that corner. Both removed
copies were unavailable to all other positions, so the eight physical
sources remain distinct. A virtual opposite source can only have been
assigned a tricoloured corner, even if the target is adjacent.
\end{proof}

\begin{lemma}[Intersecting cuts]\label{lem:intersecting-cuts}
A pairwise-intersecting family of adjacent cuts is contained, up to colour
permutation and simultaneous reversal, in one of the following families:
\[
\begin{array}{c|l}
\text{Butterfly}&C(01),C(02),C(12),C(123),C(124),C(012)\\
\text{Two stars}&C(01),C(02),C(03),C(04),C(034),C(024),C(023)\\
\text{Star and triangle}&C(01),C(02),C(03),C(04),C(034),C(024),C(014).
\end{array}
\]
\end{lemma}
\begin{proof}
Represent a size-two outset by its positive pole pair and a size-three
outset by its negative pole pair. Same-sign cuts are disjoint exactly
when their pairs are disjoint; opposite-sign cuts are disjoint exactly
when their pairs agree. The two signed edge families in $K_5$ are therefore
intersecting and share no edge. Each intersecting edge family lies in a
star or triangle: after edges $ab,ac$, an edge not containing $a$ must
be $bc$.
Two complete triangles without a common edge meet in one vertex, giving
the butterfly. For a triangle and a star, a centre outside the triangle
gives the star-and-triangle family. A centre inside permits only its two
edges to the outside vertices; extend them to a second triangle.
For two stars with different centres, assign their common edge to the
sign already using it, if either, then extend both families to obtain
the two-star form. With the same centre and at least two edges in each
family, each has exactly two edges; complete each to a triangle. If one
family has at most one edge, choose its other endpoint as centre, or any
different centre if it is empty, reducing to the different-centre case.
\end{proof}

The appendix proves that thirteen adjacent sources of class cap three in
any one of these families assemble after any two opposite sources are
adjoined. These are the butterfly, two-star and star-and-triangle assembly
arguments. The last reduces, by two small matching lemmas, to the
nine-source construction in Lemma~\ref{lem:nine-core}.

\begin{proposition}\label{prop:two-opposite-elementary}
Two opposite sources and thirteen adjacent sources of class cap three
have an assembly.
\end{proposition}
\begin{proof}
If a disjoint pair of adjacent classes occurs, compress one copy of each.
The resulting pool has three opposite and eleven adjacent sources;
Lemma~\ref{lem:elementary-inputs}(3) applies, and
Lemma~\ref{lem:compression} lifts the assembly. Otherwise
Lemma~\ref{lem:intersecting-cuts} and the three family arguments apply.
\end{proof}

\subsection{The fifteen-source conclusion}
\begin{proof}[Proof of Theorem~\ref{thm:main}]
Lemma~\ref{lem:lower} proves the lower bound. Suppose that a fifteen-source
pool fails. Any subpool fails. Five copies of one opposite variety are
excluded by Lemma~\ref{lem:elementary-concentration}, retaining them in a
fourteen-subpool. If an opposite variety occurs three or four times, the
same lemma excludes any other opposite variety.

Put $A=\sum_j\binom{g_j}{2}$, where $g_j$ are adjacent-class counts.
With four opposite copies, deleting any of the eleven adjacent sources
leaves class sizes $5,5$ and pair count $20$. Each original same-class
pair survives nine deletions, giving $9A=11\cdot20=220$, impossible.
With three opposite copies the twelve deletions instead give
$10A=12(\binom52+2\binom32)=192$, also impossible.
Every opposite variety therefore occurs at most twice. In a mixed pool,
four copies of an adjacent class could be retained with an opposite source
in a fourteen-subpool, contradicting the concentration lemma. Thus mixed
pools also have adjacent-class cap three.

Let $o$ be the number of opposite sources. For $o=0,1$, use
Lemma~\ref{lem:elementary-inputs}(1),(2). For $o=2$, use
Proposition~\ref{prop:two-opposite-elementary}. For $o=3,4$, retain three
opposite and eleven adjacent sources and use item~3. For $o=5,6$, retain
five opposite and nine adjacent sources. If the latter occupy nine
classes, item~5 applies. Otherwise some class $A$ occurs at least twice;
the cap of three ensures at least three occupied classes, so select
$2A+B+C$ and use item~4. For $o=7$, item~7 applies, and for $o\ge8$
retain eight opposite sources and use item~6. Each possibility yields an
assembly, contradicting failure. Hence $N_5=15$.
\end{proof}
